% Structural completion of constant S3 types and their isotypic shallow deformations.
\begin{lemma}[Proved: the complementary shallow cubic tangent quotient]
\label{C:shallow_tangent_quotient}
Let $V=E\oplus F\oplus Y$ over $\mathbb F_3$, with
$\dim E=\dim F=a$, $\dim Y=b$, and let
\[
 N=\begin{pmatrix}0&I_a&0\\0&0&0\\0&0&0\end{pmatrix}.
\]
Let $J$ be any $\mathbb F_3$-vector space. The tangent condition
for a cubic perturbation of a residual action $I+N$ is
$NUN=0$. Modulo infinitesimal conjugacies $U\mapsto U+[N,H]$,
each such perturbation has a unique representative
\[
 U(S,A,B,C)=
 \begin{pmatrix}0&0&0\\0&S&A\\B&0&C\end{pmatrix},
 \quad
 \begin{array}{ll}
 S\in M_a(J),& C\in M_b(J),\\
 A\in\operatorname{Hom}(Y,F)\otimes J,&
 B\in\operatorname{Hom}(E,Y)\otimes J.
 \end{array}                                                \tag{C.1.36}
\]
The image of commuting perturbations is exactly $A=B=0$.
Thus quotienting further by those perturbations leaves the
two displayed opposite arrow spaces.

Inversion of the cyclic generator, identified back with $I+N$
using $P=\operatorname{diag}(-I_a,I_a,I_b)$, acts on this
quotient presentation by
\[
          (S,A,B,C)\longmapsto(-S,-A,B,-C).              \tag{C.1.37}
\]
In particular inversion is not minus the identity on the
complementary tangent quotient. No additivity of the Swan
class in the noncommuting parameters is asserted.
\end{lemma}
\begin{proof}
The derivative of $(I+N)^3$ is
$N^2U+NUN+UN^2=NUN$. This condition says exactly that
the $F,E$ block of $U$ is zero. Direct multiplication gives
\[
 [N,H]=
 \begin{pmatrix}
 H_{FE}&H_{FF}-H_{EE}&H_{FY}\\
 0&-H_{FE}&0\\
 0&-H_{YE}&0
 \end{pmatrix}.
\]
These blocks remove $U_{EE},U_{EF},U_{EY},U_{YF}$ independently.
The surviving blocks are
$S=U_{EE}+U_{FF}$, $A=U_{FY}$, $B=U_{YE}$, and $C=U_{YY}$,
which are unchanged by every such commutator. This proves
existence and uniqueness. Commutation with $N$ says that
$U_{EE}=U_{FF}$ and $U_{FY}=U_{YE}=0$; multiplication by two
is invertible, so its image is precisely $A=B=0$.

The inverse perturbation is
$-U+NU+UN$, since $NUN=0$.
Conjugate by $P$ and then remove the three off-diagonal
gauge blocks. Its surviving blocks are exactly (C.1.37).
The calculation is linear over $\mathbb F_3$, so it also
holds for arbitrary coefficient vector spaces $J$.
\end{proof}

\begin{proposition}[Proved: the twelve constant integral $S_3$ lattices]
\label{C:cyclotomic_S3_lattices}
Let $A=\mathbb Z_3[\zeta]$ and use the nine cubic lattices of
Proposition~\ref{C:cyclotomic_C3_lattices}.
There are exactly twelve indecomposable finite $A$-free
$A[S_3]$-lattices, with unique direct-sum multiplicities:
\begin{enumerate}
\item the two signs of the natural inversion action on
$L_{\{0\}}$, $L_{\{1,2\}}$, $L_{012}$, $Q$ and $A[C_3]$;
\item $\operatorname{Ind}_{C_3}^{S_3}L_{\{1\}}$ and
$\operatorname{Ind}_{C_3}^{S_3}L_{\{0,1\}}$.
\end{enumerate}
Natural inversion swaps coordinates one and two in the
normalized character coordinates and sends $g$ to $g^{-1}$
on the regular lattice. The two induced types correspond
to the inverse-exchanged pairs
$(L_{\{1\}},L_{\{2\}})$ and
$(L_{\{0,1\}},L_{\{0,2\}})$.
The ranks are respectively $1,2,3,3,3$ for the five signed
pairs, and $2,4$ for the two induced types.
This classifies constant integral lattices, not all their
square-zero deformations.
\end{proposition}
\begin{proof}
Every one of the nine cubic indecomposables has endomorphism
ring with residue division field $\mathbb F_3$.
For $L_S$ a generic endomorphism is diagonal in its distinct
characters, and preserving the common residual line forces
its diagonal residues to agree. For $Q$, preservation of the
residual sum-zero plane imposes the same condition.
For the regular lattice the endomorphism ring is the local
order $A[C_3]$. These observations also show directly that
all these indecomposables are absolutely indecomposable in
the residue division-ring sense used here.

Restrict a proposed $S_3$-lattice to $C_3$, and write it as
$M=\bigoplus_i L_i^{m_i}$. Its endomorphism ring $E$ over
$A[C_3]$ is finite over the complete DVR $A$, hence semiperfect,
and
\[
          E/\operatorname{rad}E
             \simeq\prod_i M_{m_i}(\mathbb F_3).
\]
A transposition acts semilinearly with respect to inversion
of the cubic generator and has square one. It permutes the
factors by the five fixed types and two exchanged pairs in
the statement. In an exchanged pair the two multiplicities
must be equal, and a change of multiplicity bases makes the
action the exchange map. On a fixed type choose its stated
natural inversion $\tau_i$. Relative to $\tau_i$, the action
on the residual multiplicity space is an ordinary linear
involution: $\tau_i$ acts trivially on the residue division
field $\mathbb F_3$. Since two is invertible, that multiplicity
space splits into its plus and minus eigenspaces.

These residual normal forms determine the full lattice
action. Let $s$ be the
given transposition and let $s_0$ be the standard action
with the computed multiplicities. Lift a residual
intertwining change of basis to $P_0\in E$ and put
\[
                    P=(P_0+s_0P_0s)/2.
\]
Both summands commute with the cubic generator. The map
$P$ intertwines $s$ with $s_0$ and has the same invertible
image as $P_0$ modulo $\operatorname{rad}E$; hence $P$ is an
automorphism. This proves the asserted direct-sum description.

The residual endomorphism algebra of each listed summand
with its transposition is $\mathbb F_3$, so it is
indecomposable. The two signs on a fixed type are not
isomorphic: a residual scalar cannot intertwine $+1$ with
$-1$. The exchanged types have different cubic restrictions.
The same residual normal-form argument gives uniqueness of
all multiplicities, or equivalently one may apply
Krull--Schmidt to these finite endomorphism rings.
\end{proof}

\begin{theorem}[Proved: shallow cancellation for each constant $S_3$ isotypic type]
\label{C:S3_isotypic_shallow_cancellation}
Put $k=\mathbb Z_3[\zeta]\ltimes J$, where $J$ is any
$\mathbb F_3$-vector space. Let $M$ be a $k[S_3]$-module finite free over $k$ whose reduction $M_0=M/JM$ is a direct sum
of copies of one of the twelve indecomposable lattices in
Proposition~\ref{C:cyclotomic_S3_lattices}. Then
\[
                         [M]=[M_0\otimes_Ak]
                                  \quad\hbox{in }R_k(S_3).      \tag{C.1.38}
\]
More generally the same conclusion holds if, after removing
constant regular cubic summands, the residual cubic action
has only size-one Jordan blocks or only size-two Jordan
blocks. The theorem does not assert cancellation for arbitrary
mixtures of constant isotypic types.
\end{theorem}
\begin{proof}
Restrict first to $C_3$ and write the cubic action as $T_0+U$.
A constant regular cubic summand lifts and splits off: lift its generator to obtain a
map from $k[C_3]$ which is a split injection on underlying
free $k$-modules. Its free quotient is again a representation.
The injection splits equivariantly because $k[C_3]$ is
coinduced from the trivial subgroup. Explicitly, if $r$ is an
underlying linear retraction and $\lambda(m)$ is the identity
coefficient of $r(m)$, then
\[
       r'(m)=\sum_{g\in C_3}\lambda(g^{-1}m)g
\]
is an equivariant retraction. There is no division by three. Repeating removes all such summands.
This removal is performed only in the cubic category; no
$S_3$-equivariant splitting is needed. A regular cubic summand
and its constant reduction have trace zero on either nonidentity
cubic element, so removal preserves the relative trace difference
for each of the two generators. This argument also shows that a
deformation of a sum of regular cubic lattices has zero relative
cubic class.

For the remaining cubic types $N=\overline T_0-I$ has
$N^2=0$. The cubic tangency condition is $NUN=0$.
For the original $S_3$-representation and its constant reduction,
a three-cycle and its inverse have equal traces. Thus the two
relative trace differences agree, also after the cubic regular
summands have been removed. The inverse perturbation is $-U+NU+UN$. The block condition
$NUN=0$ gives $\operatorname{tr}(NU)=0$.
Consequently $\operatorname{tr}U=-\operatorname{tr}U$ and
$\operatorname{tr}U=0$.
If there are only size-one or only size-two residual blocks,
Lemma~\ref{C:shallow_tangent_quotient} makes every perturbation
infinitesimally conjugate to a commuting one. Its relative
cubic class is zero by Proposition~\ref{C:shallow_commuting_trace}.
To apply Lemma~\ref{C:S3_cyclic_detection}, observe that $k$
is local, so $K_0(k)=\mathbb Z\to K_0(k\otimes\mathbb Q)=\mathbb Z$
is injective. The relative class is generically zero because
$J\otimes\mathbb Q=0$. That lemma then gives (C.1.38).
This covers the character types, the signed pair-gluing
types, both exchanged-pair types, and the regular types.

It remains to treat copies of one signed $L_{012}$ or one
signed $Q$. First make the transposition equal to its
constant action $s_0$. If it is $s_0+V$, its square-one
condition says $s_0V+Vs_0=0$; conjugation by
$I-s_0V/2$ removes $V$.

For $L_{012}$, choose an integral basis with residual
$N=E_{12}$ and with third vector $Y=\pi e_0$. The natural
inversion has signs $(-,+,+)$ on the residual chain and
$Y$. Such an integral eigenbasis is explicit: take
$\pi(e_1-e_2)$, $(1,1,1)$ and $\pi e_0$.
For $Q$, take the free equivariant quotient given by the
zeroth coordinate. An integral inversion eigenbasis is
\[
 \pi(e_1+e_2),\qquad e_1-e_2,\qquad
                      e_0-(e_1+e_2)/2.
\]
It again has residual $N=E_{12}$, while inversion has
signs $(+,-,+)$; the first two vectors span the quotient's
kernel. Multiplying the whole transposition by minus one
does not alter its conjugation action on endomorphisms.
For any number of copies these sign patterns are unchanged
on the respective multiplicity blocks.

The transposition acts linearly on cubic tangent cocycles
by
\[
 \mathcal I(U)=-s_0(I-N)U(I-N)s_0.
\]
It sends a gauge commutator $[I+N,H]$ to
$[I+N,s_0Hs_0]$. The original cocycle is fixed by
$\mathcal I$ because it is an $S_3$-action.
For the $L_{012}$ sign pattern, (C.1.37) shows that a fixed
gauge class has $S=A=C=0$ and arbitrary $B$.
For the $Q$ sign pattern it has $S=B=C=0$ and arbitrary $A$.

Take the corresponding pure $B$ or pure $A$ representative
$W$ and replace it by $(W+\mathcal I(W))/2$.
This is an invariant cocycle in the same cubic gauge class.
If $U-W'$ is the commutator of a matrix $H$, averaging
$H$ with $s_0Hs_0$ makes the conjugacy preserve the
transposition as well. Thus the original representation is
isomorphic to the one with this averaged perturbation.

In the $L_{012}$ case the pure $B$ map has image in the
actual free invariant character submodule generated by the
chosen $Y$ vectors, a direct summand of the underlying free
module, and vanishes on that submodule.
The same is true after averaging, because $T_0$ and $s_0$
preserve that submodule. Hence the deformed representation
is an admissible extension of the same two constant
representations as $M_0$.
In the $Q$ case the pure $A$ map has image in the actual
free kernel of the chosen character quotient and vanishes
on that kernel. Averaging again preserves these properties,
so its subrepresentation and quotient are both constant.
The corresponding two exact sequences prove (C.1.38).

The two opposite arrow directions may both occur in mixed
constant types. The argument has not made the Swan class
additive in those directions, and therefore does not remove
that remaining gate.
\end{proof}
